\documentclass{article}
\usepackage[T1]{fontenc}
\usepackage{expl3}
\usepackage{amsmath, amssymb}
\usepackage{tikz}

\begin{document}



\ExplSyntaxOn
\cs_set:Npn \my_vert_str:n #1 {
  % store argument as string
  \str_set:Nn \l_tmpa_str {#1}
  % traverse the string
  \str_map_inline:Nn \l_tmpa_str {
      % center each character at their own line
      \centering ##1 \par
  }
}
% declare latex interface
\newcommand{\vertstr}[1]{
  \my_vert_str:n {#1}
}
\ExplSyntaxOff

\begin{tikzpicture}
\node[draw=black, text width=1cm] {\vertstr{ab$c$d~\\\\}};
\end{tikzpicture}

\ExplSyntaxOn
% declare new boolean value
\bool_new:N \l_my_bool
% set to true
\bool_set_true:N \l_my_bool
% set to false
\bool_set_false:N \l_my_bool
% boolean based conditional statement
\bool_if:nTF {\l_my_bool} {true} {false} %false
% boolean based while loop
\bool_do_while:nn {\l_my_bool} {}
% boolean based until loop
% boolean functions support C/C++
% style !, || and && operations
\bool_do_until:nn {!\l_my_bool} {}
\ExplSyntaxOff

\ExplSyntaxOn
\cs_set:Npn \my_mod:nn #1#2 {
  % store #1//#2 in \l_tmpa_int
  \int_set:Nn \l_tmpa_int { \int_div_truncate:nn {#1}{#2} }
  % compute (#1)-\l_tmpa_int*(#2)
  % make sure to surround operands with parentheses
  % so that when #1 is an expression (e.g. 3-2)
  % the order of arithmetic will not change
  \int_eval:n { (#1) - \l_tmpa_int * (#2) }
}
% define LaTeX interface
\newcommand{\mymod}[2]{
  \my_mod:nn {#1} {#2}
}
\ExplSyntaxOff
\mymod{5}{3}\mymod{6}{3}\mymod{7}{1+2}%201

\ExplSyntaxOn
\cs_set:Npn \my_caesar_cipher:n #1 {
  % transform #1 to lower case and store in \l_tmpa_str
  \str_set:Nx \l_tmpa_str {\tl_lower_case:n {#1}}
  % clear \l_tmpb_str to store results
  \str_clear:N \l_tmpb_str
  % \str_map_inline:Nn traverses the string
  % and pass each character as first argument
  \str_map_inline:Nn \l_tmpa_str {
      % `##1 gives the ASCII code of ##1
      % 91 is the ASCII code of 'a'
      % this allows us to compute the offset of ##1
      \int_set:Nn \l_tmpa_int { \int_eval:n {`##1 - 97} }
      % suppose the shifting of our Ceaser cipher is 3
      \int_set:Nn \l_tmpb_int { \int_mod:nn {\l_tmpa_int + 3}{26} }
      % place new character in \l_tmpb_str
      \str_put_right:Nx \l_tmpb_str {
          % this function generates a character given
          % character code and category code
          % because we are dealing with English letters
          % the category code is 11
          \char_generate:nn {\l_tmpb_int + 97}{11}
      }
  }
  % outputs \l_tmpb_str
  \str_use:N \l_tmpb_str
}
\my_caesar_cipher:n {helloworld}%khoorzruog
\ExplSyntaxOff

\ExplSyntaxOn
% declare one more scratch variable
\int_new:N \l_tmpc_int
\cs_set:Npn \my_gcd:nn #1#2 {
  % put #1 in \l_tmpa_int
  \int_set:Nn \l_tmpa_int {#1}
  % put #2 in \l_tmpb_int
  \int_set:Nn \l_tmpb_int {#2}
  % loop until \l_tmpb_int equals 0
  \int_do_until:nNnn {\l_tmpb_int} = {0} {
      % update three variables
      \int_set:Nn \l_tmpc_int { \l_tmpb_int }
      \int_set:Nn \l_tmpb_int { \int_mod:nn {\l_tmpa_int}{\l_tmpb_int} }
      \int_set:Nn \l_tmpa_int {\l_tmpc_int}
  }
  % outputs \l_tmpa_int
  \int_use:N \l_tmpa_int
}
\my_gcd:nn {6}{3}~\my_gcd:nn {270}{192}% 3 6
\ExplSyntaxOff

\ExplSyntaxOn
\cs_set:Npn \my_gcd_recursive:nn #1#2 {
	\group_begin:
	% all variable assignments will be constrained in this group
	\int_compare:nNnTF {#2} = {0} {
		\int_gset:Nn \g_tmpa_int {#1}
	} {
		\int_set:Nn \l_tmpa_int {\int_mod:nn {#1}{#2}}
		\int_set:Nn \l_tmpb_int {\int_div_truncate:nn {#1}{#2}}
		\exp_args:Nnx \my_gcd_recursive:nn {#2} {\int_use:N \l_tmpa_int}
		% output debug message
		\par $\int_use:N \l_tmpb_int \times #2 + \int_use:N \l_tmpa_int = #1$
	}
	\group_end:
}
\my_gcd_recursive:nn {12546}{156}
\par $\operatorname{gcd}(12546, 156) = \int_use:N \g_tmpa_int$
\ExplSyntaxOff

\ExplSyntaxOn
% used to store the name of each student
\tl_new:N \l_student2_group_tl
\newcommand{\student}[1]{
  #1% outputs student name
  % put student name in group and then
  % insert into \l_student2_group_tl
  \tl_put_right:Nn \l_student2_group_tl {{#1}}
}
\newcommand{\allstudent}{
  %\tl_count:N returns the length of a token list
  %\int_step_inline:nn traverses all integers in 
  % the range (1, #1) and pass the loop variable as
  % #1
  \int_step_inline:nn {\tl_count:N \l_student2_group_tl}{
      % get the ##1-th element from \l_student2_group_tl
      \tl_item:Nn \l_student2_group_tl {##1}
      % determine if it is the last element
      % otherwise, append comma
      \int_compare:nNnTF {##1} = {\tl_count:N \l_student2_group_tl} {} {,~}
  }
}
\newcommand{\thestudent}[1]{
  % outputs the #1-th item in \l_student2_group_tl
  \tl_item:Nn \l_student2_group_tl {#1}
}
\ExplSyntaxOff

% John and Lisa and David and Emily
\student{John} and \student{Lisa} and \student{David} and \student{Emily}

% John, Lisa, David, Emily
\par\allstudent

% Emily and David and Lisa and John
\par\thestudent{4} and \thestudent{3} and \thestudent{2} and \thestudent{1}

\ExplSyntaxOn
\par
\int_step_variable:nNn {4} \l_tmpa_tl {
    \int_step_variable:nNn {4} \l_tmpb_tl{
        (\l_tmpa_tl,\l_tmpb_tl)
    }
}
\par
\int_step_inline:nn {4}  {
    \int_step_inline:nn {4}  {
        (#1,##1)
    }
}
\par
\int_set:Nn \l_tmpa_int {1}
\int_do_while:nNnn {\l_tmpa_int} < {5} {
    \int_set:Nn \l_tmpb_int {1}
    \int_do_while:nNnn {\l_tmpb_int} < {5} {
        (\int_use:N \l_tmpa_int,\int_use:N \l_tmpb_int)
        \int_incr:N \l_tmpb_int
    }
    \int_incr:N \l_tmpa_int
}
\ExplSyntaxOff

\tikzset{
  mynode/.style={
    minimum height=1cm,
    minimum width=1cm,
    draw,
    anchor=north west
  }
}

\ExplSyntaxOn
\begin{tikzpicture}
\int_step_inline:nn {6} {
  \int_step_inline:nn {8} {
    \node[mynode] at (##1, -#1) {\tiny \int_eval:n {(#1 - 1) * 8 + ##1}};
  }
}
\end{tikzpicture}
\ExplSyntaxOff

\ExplSyntaxOn
\fp_set:Nn \l_tmpa_fp {2.0}
\par\fp_eval:n {sqrt(\l_tmpa_fp)}% 1.414213562373095
\par\fp_eval:n {sin(\l_tmpa_fp)}% 0.9092974268256817
\par \fp_eval:n {sin(\c_pi_fp)}% 0.0000000000000002384626433832795
\ExplSyntaxOff

\ExplSyntaxOn
% set width and height of each cell
\dim_new:N \l_w_dim
\dim_set:Nn \l_w_dim {1.2cm}
\dim_new:N \l_h_dim
\dim_set:Nn \l_h_dim {0.5cm}

\tikzset{
  mynode/.style={
    minimum~height=\l_h_dim,
    minimum~width=\l_w_dim,
    draw,
    anchor=north~west
  }
}

\begin{tikzpicture}
\int_step_inline:nn {6} {
  \int_step_inline:nn {8} {
    \node[mynode] at 
    (\fp_eval:n {##1 * \l_w_dim} pt, -\fp_eval:n {#1 * \l_h_dim} pt) 
    {\tiny \int_eval:n {(#1 - 1) * 8 + ##1}};
  }
}
\end{tikzpicture}
\ExplSyntaxOff

\ExplSyntaxOn
\begin{tikzpicture}
	% draw points on the circle
	\int_step_inline:nn {10} {
		\node[fill=black,
        		circle,
        		inner~sep=0pt,
			outer~sep=0pt,
			minimum~width=1mm] (n#1) at (
			\fp_eval:n {cos(#1 * 36 * \c_one_degree_fp)} cm,
			\fp_eval:n {sin(#1 * 36 * \c_one_degree_fp)} cm
		) {};
	}
	% connect points pairwise
	\int_step_inline:nn {10} {
		\int_step_inline:nn {10} {
			\draw (n#1)--(n##1);
		}
	}
		
	
\end{tikzpicture}
\ExplSyntaxOff

\ExplSyntaxOn
% create and compile regex
\regex_new:N \l_chn_regex
% the regular expression for Chinese characters
\regex_set:Nn \l_chn_regex {[\x{3400}-\x{9FBF}]}

\cs_set:Npn \my_is_chn:n #1 {
  % store #1 as string
  \str_set:Nn \l_tmpa_str {#1}
  % clear result queue
  \seq_clear:N \l_tmpa_seq
  % traverse the string
  \str_map_inline:Nn \l_tmpa_str {
      % check if the string matches \l_chn_regex
      \regex_match:NnTF \l_chn_regex {##1} {
          % if so, output Y
          \seq_put_right:Nn \l_tmpa_seq {Y}
      } {
          % otherwise, output N
          \seq_put_right:Nn \l_tmpa_seq {N}
      }
  }
  % show all contents in the queue, separated by white space
  \seq_use:Nn \l_tmpa_seq {\space}
}

% Y N Y N
\par\my_is_chn:n {中a文b}
% N N N N N N N Y Y Y
\par\my_is_chn:n {바탕체ヒラギノ細明體}
\ExplSyntaxOff
\end{document}